\section{补充统计与图件}

\begin{table}[!htbp]
\centering\small
\caption{\textbf{输入变量与派生指标的定义。} SCF、SWE 和 T2m 均来自提供的月度数据，按每月实际天数聚合为 DJF 冬季值；面积权重和固定分析掩膜来自同一数据包。}\label{tab:variables}
\setlength{\tabcolsep}{5pt}
\begin{tabular}{lll}
\toprule
指标 & 单位 & 定义及统计口径 \\
\midrule
SCF & 比例（0--1） & 格点积雪覆盖比例 \\
SWE & mm & 格点雪水当量；空间代表性定义未说明 \\
T2m & $^\circ$C & 格点 2 米气温 \\
SCA & km$^2$ & 固定域内面积乘 SCF 后求和 \\
SCR & \% & SCA 除以固定域总面积 \\
$V_A$ & km$^3$ & SWE 为全格点平均时的储量估计 \\
$V_B$ & km$^3$ & SWE 为有雪面积条件平均时的储量估计 \\
\bottomrule
\end{tabular}
\end{table}

\begin{table}[!htbp]
\centering\footnotesize
\caption{\textbf{不同 SCF 阈值下的早晚期边界配对位移。} 1980--1989 年与 2013--2022 年各含 10 个 DJF 冬季。只在两期均可识别交叉的经度上配对，报告位移中位数及 10 经度空间区块自助法 95\% CI。正值表示北移；区域为近似经纬度窗口。}\label{tab:boundary}
\setlength{\tabcolsep}{5pt}
\begin{tabular}{llrr}
\toprule
SCF 阈值 & 范围 & 配对经度数 & 纬度位移 ($^\circ$) [95\% CI] \\
\midrule
0.25 & 全分析域 & 169 & $0.275\;[0.117,\ 0.518]$ \\
0.25 & 欧洲窗口 & 46 & $0.177\;[0.088,\ 0.328]$ \\
0.25 & 亚洲窗口 & 68 & $0.195\;[0.049,\ 0.345]$ \\
0.25 & 北美洲窗口 & 55 & $0.516\;[0.103,\ 1.716]$ \\
0.50 & 全分析域 & 174 & $0.366\;[0.207,\ 0.593]$ \\
0.50 & 欧洲窗口 & 46 & $0.450\;[0.168,\ 0.863]$ \\
0.50 & 亚洲窗口 & 68 & $0.419\;[0.103,\ 1.290]$ \\
0.50 & 北美洲窗口 & 60 & $0.263\;[0.173,\ 0.553]$ \\
0.75 & 全分析域 & 185 & $0.220\;[0.080,\ 0.380]$ \\
0.75 & 欧洲窗口 & 45 & $0.287\;[0.173,\ 1.390]$ \\
0.75 & 亚洲窗口 & 71 & $0.527\;[0.140,\ 0.792]$ \\
0.75 & 北美洲窗口 & 68 & $0.025\;[0.000,\ 0.220]$ \\
\bottomrule
\end{tabular}

\end{table}

\begin{figure}[!htbp]
\centering
\includegraphics[width=\linewidth]{figures/fig07b_snowline_threshold_sensitivity.pdf}
\caption{\textbf{阈值 0.25 和 0.75 的积雪覆盖边界。} 虚线和实线分别表示 1980--1989 年与 2013--2022 年各 10 个 DJF 冬季的平均 SCF 等值线，底图为晚减早 SCF（pp）。}\label{fig:thresholdsens}
\end{figure}
